% Copyright (c) 2025 ZKLib Contributors. All rights reserved.
% Released under Apache 2.0 license as described in the file LICENSE.
% Authors: Poulami Das (Least Authority)

\section{STIR}\label{sec:stir}

STIR is a Reed--Solomon IOPP whose soundness combines quotienting, out-of-domain
sampling, folding, and a random linear combination proximity generator.  The current ArkLib
formalization has two layers.  The paper-level theorem statements
\texttt{StirIOP.stir\_main} and \texttt{StirIOP.stir\_rbr\_soundness} are front doors.  The
protocol layer builds single-round, multi-round, and checking-verifier objects and proves their
completeness.  The remaining soundness work is isolated as named correlated-agreement and
round-by-round bridge residuals.  The live implementation tracker is
\href{https://github.com/lalalune/ArkLib/issues/301}{STIR issue \#301}; this blueprint records
that tracker's current Lean state and should not be read as unconditional closure of the paper
soundness theorem.

\subsection{Tools for Reed--Solomon codes}\label{sec:stir_tools}

\subsubsection{Random linear combination as a proximity generator}\label{sec:proximity_gap}

\begin{theorem}[STIR proximity gap front door]\label{thm:proximity_gap}
\lean{STIR.proximity_gap,STIR.proximity_gap_iff_jointAgreement}
\uses{def:distance_from_code,def:reed_solomon_code,def:list_close_codewords,def:correlated_agreement}
    Let $\code := \rscode[\field,\evaldomain,\degree]$ and let
    $B^*(\rate)=\sqrt{\rate}$.  For $\distance < 1-B^*(\rate)$, if the random
    power-combination $\sum_j r^j f_j$ is close to $\code$ with probability above the
    BCIKS/STIR error threshold, then the words $f_j$ have a common large agreement set with
    Reed--Solomon codewords.  The Lean statement explicitly constrains the generator by
    \texttt{hGen : GenFun r j = r \^{} j}; without this power-generator condition the older
    arbitrary-generator statement is false.  The theorem is also bridged to the shared
    \texttt{Code.jointAgreement} API.
\end{theorem}

\begin{remark}[Proximity-gap status]\label{rem:stir_proximity_gap_status}
    The STIR proximity-gap theorem is a front-door predicate.  The active proof route goes
    through the BCIKS20 polynomial-curve correlated-agreement machinery, whose Johnson-regime
    extraction is exposed through named residuals such as
    \texttt{ProximityGap.StrictCoeffPolysResidual}.  Small-field and vacuous routes discharge
    parts of this family in dedicated modules.
\end{remark}

\subsubsection{Univariate function quotienting}\label{sec:quotienting}

\begin{definition}[Function quotient]\label{def:quotient}
\lean{Quotienting.funcQuotient}
    For a word $f$, a finite answer set $S$, answer values $\mathsf{Ans}$, and fill values
    $\mathsf{Fill}$, \texttt{funcQuotient} returns $\mathsf{Fill}(x)$ on $S$ and
    $(f(x)-\widehat{\mathsf{Ans}}(x))/V_S(x)$ away from $S$, where $V_S$ is the vanishing
    polynomial of $S$.
\end{definition}

\begin{definition}[Polynomial quotient]\label{def:poly_quotient}
\lean{Quotienting.polyQuotient}
\uses{def:quotient}
    The polynomial quotient is $(\hat f-\widehat{\mathsf{Ans}})/V_S$, with
    $\widehat{\mathsf{Ans}}$ interpolating $\hat f$ on $S$.
\end{definition}

\begin{lemma}[Quotienting lemma]\label{lemma:quotienting}
\lean{Quotienting.quotienting}
\uses{def:quotient,def:poly_quotient,def:reed_solomon_code,def:distance_from_code,def:list_close_codewords}
    If every $\distance$-close Reed--Solomon codeword disagrees with the supplied answer on
    some point of $S$, then the quotient function is far from the smaller-degree
    Reed--Solomon code, up to the explicit disagreement-set correction.  The Lean theorem
    includes the domain-identity and degree-fit hypotheses needed for interpolation
    uniqueness.
\end{lemma}

\subsubsection{Out-of-domain sampling}\label{sec:out_of_domain_smpl}

\begin{lemma}[Out-of-domain collision bound]\label{lemma:out_of_domain_smpl}
\lean{OutOfDomSmpl.out_of_dom_smpl_1,OutOfDomSmpl.out_of_dom_smpl_2}
\uses{def:reed_solomon_code,def:list_decodable,def:list_close_codewords}
    If a Reed--Solomon code is list decodable, then repeatedly sampling points outside the
    evaluation domain bounds the probability that two distinct nearby codewords agree on all
    sampled points.  The formal proof reduces the probability bound to finite counting over the
    domain complement.
\end{lemma}

\subsubsection{Folding univariate functions}\label{sec:folding_uf}

\begin{lemma}[Polynomial folding decomposition]\label{fact:poly_folding}
\lean{Polynomial.FoldingPolynomial.folding_polynomial_is_unique,Polynomial.FoldingPolynomial.substitution_property_of_folding_polynomial,Polynomial.FoldingPolynomial.folding_polynomial_deg_y_bound,Polynomial.FoldingPolynomial.folding_polynomial_deg_x_bound}
    A univariate polynomial can be uniquely represented as a bivariate polynomial in
    $(q(Z),Z)$ with the expected degree bounds.  This is the algebraic basis for STIR folding.
\end{lemma}

\begin{definition}[Polynomial folding]\label{def:poly_folding}
\lean{Polynomial.FoldingPolynomial.polyFold}
\uses{fact:poly_folding}
    Given $\hat f$, a folding parameter $k$, and a challenge $r$, \texttt{polyFold} evaluates the
    bivariate decomposition of $\hat f$ at $Y=r$.
\end{definition}

\begin{definition}[Function folding]\label{def:fn_folding}
\lean{ProximityGap.foldWord,ProximityGap.foldValue}
\uses{def:poly_folding}
    The function fold interpolates the values on a $k$-fiber and evaluates the resulting
    degree-$<k$ polynomial at the verifier challenge.
\end{definition}

\begin{lemma}[Folding preserves distance]\label{lemma:folding}
\lean{ProximityGap.folding_preserves_distance}
\uses{def:reed_solomon_code,def:fn_folding,def:distance_from_code,thm:proximity_gap}
    If a word is far from the Reed--Solomon code, then with probability controlled by the
    STIR proximity-gap error its fold remains far from the folded Reed--Solomon code.
\end{lemma}

\subsubsection{Combining functions of varying degrees}\label{sec:combine}

\begin{lemma}[Geometric sum]\label{fact:geometric_sum}
\lean{Combine.geometric_sum_units}
    The finite geometric sum is expressed in closed form when the ratio is not $1$, and as
    $a+1$ when the ratio is $1$.
\end{lemma}

\begin{definition}[Degree-correcting combination]\label{def:combine}
\lean{Combine.combine,Combine.degCor}
\uses{fact:geometric_sum}
    \texttt{Combine.combine} forms the STIR random combination of functions with different
    target degrees by multiplying each function by the appropriate geometric degree-correction
    factor.  \texttt{degCor} is the one-function specialization.
\end{definition}

\begin{lemma}[Combination proximity theorem]\label{lemma:combine}
\lean{Combine.combine_theorem,Combine.combine_theorem_proximityError,Combine.combine_theorem_of_card_le,Combine.combine_theorem_of_card_le_e7,Combine.combine_theorem_proximityError_of_card_le}
\uses{def:reed_solomon_code,def:combine,thm:proximity_gap}
    If the degree-corrected random combination is close to the target Reed--Solomon code with
    probability above the proximity-gap threshold, then all input functions agree with
    codewords on a common large set.  The main theorem consumes the BCIKS strict-coefficient
    residual explicitly; the \texttt{card\_le} and \texttt{card\_le\_e7} variants discharge it
    in vacuous/small-field regimes.
\end{lemma}

\subsection{STIR front-door statements}\label{sec:stir_main_theorems}

\begin{definition}[STIR parameters and relation]\label{def:stir_params}
\lean{StirIOP.Params,StirIOP.ParamConditions,StirIOP.Distances,StirIOP.CodeParams,StirIOP.stirRelation}
\uses{def:reed_solomon_code,def:list_decodable,def:distance_from_code}
    The Lean parameter structures store the initial degree, per-round folding parameters,
    evaluation embeddings, repetition parameters, distance targets, and list-decoding
    assumptions.  The relation \texttt{stirRelation} states relative proximity to the
    round-zero Reed--Solomon code.
\end{definition}

\begin{theorem}[STIR main theorem]\label{thm:stir}
\lean{StirIOP.stir_main}
\uses{def:stir_params,lemma:rnd_by_rnd_soundness}
    The paper-level main theorem asserts existence of a vector IOPP with the stated
    round-by-round soundness, round complexity, proof length, and query complexity bounds.
    In Lean this is a \texttt{Prop} front door; downstream construction theorems instantiate
    the existential under explicit security and complexity hypotheses.
\end{theorem}

\begin{lemma}[STIR round-by-round soundness]\label{lemma:rnd_by_rnd_soundness}
\lean{StirIOP.stir_rbr_soundness}
\uses{def:stir_params,lemma:folding,lemma:out_of_domain_smpl,lemma:quotienting,lemma:combine}
    Lemma 5.4 packages the STIR per-round error family
    $(\epsilon^{\mathsf{fold}},\epsilon^{\mathsf{out}}_i,
    \epsilon^{\mathsf{shift}}_i,\epsilon^{\mathsf{fin}})$ and asserts existence of a secure
    vector IOPP with challenge count $2M+2$.  The current formalization provides several
    conditional front doors that instantiate the vector IOPP and leave only the honest
    soundness bridge and numeric free-parameter inequalities as hypotheses.
\end{lemma}

\subsection{Protocol assembly}\label{sec:stir_protocol_assembly}

\begin{definition}[Single-round and block reductions]\label{def:stir_round_reductions}
\lean{StirIOP.Round.stirRoundReduction,StirIOP.Round3.stirRound3Reduction,StirIOP.Round3.stirFullReduction}
\uses{def:combine,def:stir_params}
    The implementation starts from a single STIR round, groups the paper's round shape into
    three-message blocks, and assembles block chains into a full reduction.
\end{definition}

\begin{theorem}[Round completeness]\label{thm:stir_round_completeness}
\lean{StirIOP.Round.stirRoundReduction_perfectCompleteness,StirIOP.Round.stirRoundVectorReduction_perfectCompleteness}
\uses{def:stir_round_reductions}
    The single-round oracle reduction and its vectorized version are perfectly complete.
\end{theorem}

\begin{definition}[Multi-round vector IOP]\label{def:stir_multi_round_iop}
\lean{StirIOP.MultiRound.stirMultiVSpec,StirIOP.MultiRound.stirMultiRoundProver,StirIOP.MultiRound.stirMultiRoundVerifier,StirIOP.MultiRound.stirMultiRoundIOP}
\uses{def:stir_params}
    The multi-round wire shape \texttt{stirMultiVSpec} has $3M+3$ protocol rounds and $2M+2$
    verifier challenges.  The assembled shell IOP uses the honest folding prover and a
    forwarding verifier.
\end{definition}

\begin{theorem}[Multi-round assembly front doors]\label{thm:stir_multi_round_front_doors}
\lean{StirIOP.MultiRound.stirMultiRoundIOP_perfectCompleteness,StirIOP.stir_rbr_soundness_of_secure_vectorIOP,StirIOP.stir_rbr_soundness_of_residuals,StirIOP.stir_main_of_secure_vectorIOP,StirIOP.stir_main_of_residuals}
\uses{def:stir_multi_round_iop,thm:stir,lemma:rnd_by_rnd_soundness}
    The shell multi-round IOP is perfectly complete, and any secure-with-gap vector IOP over
    the same wire shape instantiates the paper front doors.  The shell-specific residual is
    intentionally kept as a hypothesis rather than treated as a genuine small-error soundness
    proof.
\end{theorem}

\subsection{Checking verifier and residuals}\label{sec:stir_checking_verifier}

\begin{definition}[Checking verifier]\label{def:stir_checking_verifier}
\lean{StirIOP.MultiRound.checkingComp,StirIOP.MultiRound.checkingBool,StirIOP.MultiRound.stirCheckingVerifier,StirIOP.MultiRound.stirCheckingIOP}
\uses{def:stir_multi_round_iop,def:reed_solomon_code}
    The checking verifier performs a round-zero fold check, out/shift consistency checks between
    consecutive folded oracles, and a final low-degree membership check.  This replaces the shell
    verifier by a verifier whose acceptance predicate constrains the transcript.
\end{definition}

\begin{lemma}[Checking predicate decomposition]\label{lem:stir_checking_predicate}
\lean{StirIOP.MultiRound.checkingBool_eq_true_iff,StirIOP.MultiRound.checkingVerifier_acceptance_iff_checkingBool,StirIOP.MultiRound.checkingVerifier_acceptance_implies_fold_check,StirIOP.MultiRound.checkingVerifier_acceptance_implies_round_consistency,StirIOP.MultiRound.checkingVerifier_acceptance_implies_final_in_code}
\uses{def:stir_checking_verifier}
    Acceptance of the checking verifier is equivalent to the fold check, every out/shift
    adjacent-round consistency check, and final Reed--Solomon membership.  These facts are the
    local verifier facts required by the future probabilistic soundness bridge.
\end{lemma}

\begin{theorem}[Checking IOP completeness]\label{thm:stir_checking_completeness}
\lean{StirIOP.MultiRound.stirMultiRoundProver_runToRound_invariant,StirIOP.MultiRound.checkingBool_honest,StirIOP.MultiRound.stirCheckingIOP_perfectCompleteness}
\uses{def:stir_checking_verifier,thm:stir_round_completeness}
    The honest multi-round prover maintains the expected transcript invariant, the checking
    predicate accepts honest codewords, and the checking IOP is perfectly complete.
\end{theorem}

\begin{definition}[Checking soundness residuals]\label{def:stir_checking_residuals}
\lean{StirIOP.MultiRound.stirCheckingRbrSoundnessResidual,StirIOP.MultiRound.stirCheckingCABridge,ArkLib.ProofSystem.Stir.ErrorAccumulation.PerRoundProximityGap}
\uses{def:stir_checking_verifier,thm:proximity_gap}
    \texttt{stirCheckingRbrSoundnessResidual} is the genuine round-by-round knowledge-soundness
    obligation for the checking verifier.  \texttt{stirCheckingCABridge} states that this
    obligation follows from the BCIKS/STIR correlated-agreement residual family plus the
    per-round proximity-gap accounting keystone.
\end{definition}

\begin{theorem}[Checking IOP front doors]\label{thm:stir_checking_front_doors}
\lean{StirIOP.stir_rbr_soundness_of_checkingIOP_CA,StirIOP.stir_rbr_soundness_of_checkingIOP_card_le,StirIOP.stir_rbr_soundness_of_checkingIOP_card_le_e7,StirIOP.stir_rbr_soundness_of_checkingIOP_large,StirIOP.stir_main_of_checkingIOP_CA,StirIOP.stir_main_of_checkingIOP_card_le,StirIOP.stir_main_of_checkingIOP_card_le_e7,StirIOP.stir_main_of_checkingIOP_large}
\uses{thm:stir,lemma:rnd_by_rnd_soundness,def:stir_checking_residuals,thm:stir_checking_completeness}
    The headline STIR front doors are instantiated with the checking IOP.  The \texttt{CA}
    versions consume the correlated-agreement and bridge residuals; the \texttt{card\_le},
    \texttt{card\_le\_e7}, and \texttt{large} versions discharge the strict-coefficient
    residual family in the corresponding BCIKS routes and leave only the protocol-level bridge.
\end{theorem}


\subsection{Genuine sub-unit soundness, tightness, and repetition}\label{sec:stir_subunit}

\begin{theorem}[Multi-flip knowledge state functions]\label{thm:stir_multiflip_ksf}
\lean{ThresholdKSF.predKSF,ThresholdKSF.rbrKnowledgeSoundness_of_flipBounds}
\uses{def:stir_checking_residuals}
    An un-truncated round-indexed predicate family that matches the input relation on the
    empty transcript, never flips upward at prover rounds, and is forced at the final round
    by acceptance, is a knowledge state function; round-by-round knowledge soundness follows
    from per-challenge bounds on the genuine flip events
    $\neg\mathrm{pred}(\text{before}) \wedge \mathrm{pred}(\text{after the fresh challenge})$.
    This generalizes the single-threshold state functions used by the WHIR discharge and is
    the upgrade required to consume challenge-dependent flip estimates at sub-unit budgets.
\end{theorem}

\begin{theorem}[Genuine budget discharge]\label{thm:stir_subunit_discharge}
\lean{StirIOP.MultiRound.stirCheckingPred,StirIOP.MultiRound.stirEpsStar,StirIOP.MultiRound.stirCheckingRbrSoundness_genuine,StirIOP.MultiRound.stirCheckingCABridge_genuine,StirIOP.MultiRound.stirCheckingIOP_isSecureWithGap_genuine,StirIOP.MultiRound.stir_main_of_checkingIOP_genuine}
\uses{thm:stir_multiflip_ksf,lem:stir_checking_predicate,thm:stir_checking_completeness}
    The retired-prefix winnable state predicate (committed point-checks pass, the
    fully-committed pair with pending binding challenge is locked, and the last committed
    message is a codeword) yields round-by-round knowledge soundness of the checking
    verifier \emph{outright} at the explicit budget $\varepsilon^\star$: zero at the fold and
    shift challenges, $(|F|-(\lfloor\delta|\iota|\rfloor+1))/|F|$ at the round-two input-link
    challenge, and $(|F|-1)/|F|$ at the later pair-binding challenges.  In the single-domain
    identity-fold wire model no correlated-agreement input is needed, so the named bridge
    residual holds a fortiori, the checking IOP is secure-with-gap with no residual
    hypotheses, and Theorem~\ref{thm:stir} is instantiated with a proven soundness leg.
\end{theorem}

\begin{theorem}[Switch-prover tightness]\label{thm:stir_switch_attack}
\lean{StirIOP.MultiRound.stirChecking_switch_decision,StirIOP.MultiRound.stirChecking_switch_attack}
\uses{thm:stir_subunit_discharge}
    Against the constant nearest-codeword strategy the checking decision reduces to the
    single input-link check, and acceptance has probability at least
    $|\{x \mid f(x)=u(x)\}|/|F|$ over the round-two challenge.  Hence no budget family
    summing below $(|\iota|-\Delta(f,\mathrm{code}))/|F|$ is sound for this verifier:
    $\varepsilon^\star$ is essentially tight, and $2^{-\lambda}$ budgets at large $\lambda$
    are unachievable in the single-query wire model.
\end{theorem}

\begin{theorem}[The $t$-repetition wire model]\label{thm:stir_rep_wire}
\lean{StirIOP.MultiRound.stirMultiVSpecRep,StirIOP.MultiRound.stirCheckingIOPRep_perfectCompleteness,StirIOP.MultiRound.stirEpsStarRep,StirIOP.MultiRound.stirCheckingRepRbrSoundness_genuine,StirIOP.MultiRound.stirCheckingIOPRep_isSecureWithGap_genuine,StirIOP.MultiRound.stir_main_of_checkingIOP_rep}
\uses{thm:stir_subunit_discharge,thm:stir_switch_attack}
    Repeating every binding check at all $t$ coordinates of one length-$t$ vector challenge
    preserves perfect completeness and yields round-by-round knowledge soundness at the
    $t$-fold power budgets $((|F|-(\lfloor\delta|\iota|\rfloor+1))/|F|)^t$ and
    $((|F|-1)/|F|)^t$, via a product-marginal domination lemma for one uniformly drawn
    vector challenge.  Every security-parameter target is met at sufficiently large $t$,
    giving the first fully proven soundness leg for Theorem~\ref{thm:stir} at arbitrary
    security parameters within the identity-fold model.
\end{theorem}

\begin{remark}[Current STIR status]\label{rem:stir_status}
    Protocol objects, completeness, and the soundness layer are represented in Lean
    end-to-end: the checking verifier's round-by-round knowledge soundness is proven
    outright at its genuine budget (Theorem~\ref{thm:stir_subunit_discharge}), shown
    essentially tight (Theorem~\ref{thm:stir_switch_attack}), and lifted to arbitrary
    security parameters by $t$-fold repetition (Theorem~\ref{thm:stir_rep_wire}).  The
    checker remains a single-domain model: its folds are identity folds, so paper-STIR's
    degree-reduction behavior and the Johnson-regime correlated-agreement budgets are
    represented by the conditional front doors of
    Theorem~\ref{thm:stir_checking_front_doors}, whose remaining inputs are tracked by the
    BCIKS correlated-agreement issues.
\end{remark}
